% -------------------------------------------------------------------------------------
\section{为什么是 SU(3) 规范理论}\label{sec:qcd-why}

20 世纪 60 年代初，已知的强子已有上百种，它们的质量与量子数呈现出近似的 SU(3) 味对称性：自旋 $\frac12$ 的重子组成八重态，自旋 $\frac32$ 的重子组成十重态，赝标量介子组成八重态。Gell-Mann\cite{Gell-Mann:1964ewy}与 Zweig\cite{Zweig:1964jf}指出，这些多重态正是三种自旋 $\frac12$ 的基本组分（今天称为 $u$、$d$、$s$ 夸克）的 SU(3) 张量积分解：重子由三个夸克组成，介子由一个夸克与一个反夸克组成。

\paragraph{颜色。}夸克模型立刻遇到一个统计学上的矛盾。十重态中的 $\Delta^{++}$ 由三个 $u$ 夸克组成；自旋 $\frac32$ 要求三个夸克的自旋平行，自旋波函数对称；作为同类中最轻的态，它的空间波函数也应当对称。于是三个全同费米子的波函数在交换下对称，违背 Pauli 原理。Greenberg\cite{Greenberg:1964pe}与 Han、Nambu\cite{Han:1965pf}的解决办法形式各异，共同的内容是：夸克还带有一个取三个值的新自由度，今天称为颜色；重子的波函数在颜色指标上完全反对称，$\epsilon_{ijk}q^iq^jq^k$，从而整体反对称。所有观测到的强子都是颜色单态。颜色的个数 $N_c$ 可以直接计数：$e^+e^-$ 湮灭为强子的截面与 $e^+e^-\to\mu^+\mu^-$ 的截面之比在最低阶等于
\begin{equation}\label{eq:qcd-R0}
R_0=N_c\sum_qe_q^2 ,
\end{equation}
求和遍及质心能量以下的夸克味道（第~6~章的\cref{cor:ee-R-alphas}；$e_q$ 是以质子电荷为单位的夸克电荷）。对 $u,d,s$ 它是 $\frac23N_c$，加上 $c$ 是 $\frac{10}9N_c$，再加上 $b$ 是 $\frac{11}9N_c$；$N_c=3$ 时依次为 $2$、$\frac{10}3$、$\frac{11}3$。$R$ 的测量因而是颜色数的直接检验，其 $\order{\as}$ 修正在\cref{ex:qcd-ee-real-virtual}中计算。

\paragraph{短距离上的近自由。}同一时期，深度非弹散射显示出 Bjorken 预言的标度律\cite{Bjorken1969}：结构函数近似只依赖于无量纲变量 $x=Q^2/(2P\cdot q)$，而不单独依赖于动量转移 $Q^2$。这表明在短距离上强子的行为像是由几乎自由的点状组分构成，即 Feynman 的部分子图像\cite{Feynman1969}（第~7~章）。于是描述强相互作用的理论必须在短距离上变弱。Coleman 与 Gross\cite{Coleman:1973sx}证明了（在单圈层次上），四维可重整理论中只有非阿贝尔规范理论才具有这种渐近自由；Gross 与 Wilczek\cite{GrossWilczek1973a}、Politzer\cite{Politzer1973}计算了非阿贝尔规范理论的 $\beta$ 函数，发现夸克味道不太多时耦合随能标升高而减小（\cref{sec:qcd-af}）。把颜色 SU(3) 规范化，规范玻色子（胶子）组成色八重态\cite{Fritzsch:1973pi}：这就是 QCD。

\paragraph{什么被证明了，什么没有。}QCD 的微扰论是一个明确定义的数学对象：正规化并重整化之后，每个 Green 函数都是耦合常数的形式幂级数，每个系数是有限个 Feynman 积分之和。关于这些系数的陈述——规范参数无关性、可重整性、Slavnov--Taylor 恒等式、渐近自由、红外安全量的逐阶有限性——是可以严格证明的定理（其中一部分的完整证明很长，本书给出结构与文献）。与此相对，QCD 作为一个非微扰的四维量子场论是否存在，至今没有被构造出来：四维 Yang--Mills 理论的存在性与质量间隙是一个公开的数学问题。本书把它列为第~1~章的全局假设\cref{asm:intro-wightman}（\tagC），格点规范理论\cite{Wilson:1974sk}的数值结果是它的非微扰证据而不是证明；夸克禁闭同样没有从第一原理推导出来，不在本书范围之内。因此，凡涉及强子的陈述，最多只能达到“有严格的微扰论内容加编号假设”的地位。

\paragraph{如何阅读本章的标签。}本书的每个编号陈述都带有严格性标签（定义见第~1~章的\cref{def:intro-tags}）：\tagR\ 严格，\tagA\ 全阶论证，\tagP\ 部分，\tagC\ 猜想，\tagV\ 破坏。本章只出现 \tagR，并且总带一个说明严格性来源的限定语，作为陈述的第一句：$R_{\mathrm{def}}$ 表示结论由场与算符的定义及代数直接推出，例如规范协变性与 BRST 变换的幂零性；$R_{\mathrm{PT}}$ 表示结论在重整化微扰论中对红外正规化的部分子态的每一阶都成立，例如可重整性与渐近自由。定义与纯数学引理不带标签，例与注也不带标签。每个陈述的最后一行“依赖：”列出它的定义与证明中用到的全部更早的陈述，这些记录仿照 Lean 形式化中的蓝图，在附录~P 中自动汇成依赖图。

